\section{From frozen confidence to current-curvature scores}\label{sec:abstract-transport}

Consider a learner choosing from a finite list of available actions $\cA_t$
at round $t$.  After observing the context, it predicts each action's mean
reward $\mu_t^*(a)$ and chooses an action $a_t$ whose reward could plausibly
be high.  Only the chosen reward is then revealed.
An \emph{upper confidence bound} (UCB) makes this rule precise: its score adds
an uncertainty allowance to a prediction.  We call the score \emph{optimistic}
when it upper-bounds the true mean reward of every available action on the
stated confidence and certificate event.  This is a claim about mean rewards,
not about every noisy reward that the learner might observe.

\paragraph{What changes when the model is retrained?}
The feature vector used to interpret a past reward was fixed before that
reward arrived.  Keeping these collection-time vectors gives a frozen
reference matrix with the sequential update used by the confidence argument.
Recomputing the past vectors at the current model can give a different matrix,
even though no new rewards have been collected.  The question is how to use
this current matrix without silently assuming that the original confidence
argument still applies.

The prediction query $\bq_t(a)$ is the vector whose uncertainty matters for
action $a$; in the nonlinear specialization it is the current prediction
gradient.  A metric here is a symmetric positive-definite (SPD) matrix: its
quadratic form is positive in every nonzero direction.  A query's \emph{width}
is its norm under the inverse metric.  It measures a geometric scale, not a
confidence level by itself.  Statistical confidence also needs a valid
nonnegative radius $\beta_t$, a prediction center $m_t(a)$, and a nonnegative
error allowance $b_t(a)$.  A smaller width after retraining therefore does not, by itself,
show that prediction error has become smaller.  A \emph{certificate} is a
bound that guarantees a required comparison or numerical property; measuring
a small discrepancy in an experiment is not itself such a guarantee.

\paragraph{Start with an exact current matrix and exact widths.}
For the moment, suppose the learner computes the current matrix and its
width $s_t(a)$ exactly.  Suppose also that a statistical argument bounds the
absolute prediction error by $\beta_t\bar s_t(a)+b_t(a)$ in the frozen matrix.
A valid path certificate $\bar D_t$ supplies the comparison
$\bar s_t(a)\le e^{\bar D_t/2}s_t(a)$ for every action.  Thus the exact-case
score is
\[
 U_t(a)=m_t(a)+\beta_t e^{\bar D_t/2}s_t(a)+b_t(a).
\]
On the joint confidence and certificate event, this score upper-bounds the
mean reward.  The multiplier enlarges the current width enough to retain
the supplied reference confidence.  The center must be the one justified by
that confidence argument; it need not be the raw nonlinear prediction.
\Cref{sec:corrected-center} shows how to construct it.

This exact case explains what transport provides: a justification for a score
built around current curvature.  Direct reference scoring remains an
alternative.  The transport argument alone establishes neither a smaller
uncertainty allowance nor better action choices.  The general result next
also allows an approximate matrix, numerical width computation, and an
approximately maximizing action selector.

\paragraph{How to read the argument.}
We first assume a confidence band in the frozen metric and state the
conditions that let a current-metric score use it.  The main result is
\Cref{thm:transported-ucb}; its interpretation follows the statement.
\Cref{sec:logarithmic-transport} then supplies geometric comparisons between
metrics, and \Cref{sec:corrected-center} constructs the confidence band for
squared loss.  \Cref{sec:closed-rates-new} combines these pieces into regret
rates and explains the sharper comparison available for scaled tanh.
The finite-action picture above is enough to follow these constructions;
the general protocol below also covers action spaces that cannot be listed.

\paragraph{The matrices and their widths.}
Let $\bq_t(a)\in\R^d$ be the prediction query.  We distinguish the frozen
reference matrix $\bVbar_t$, the exact current matrix $\bV_t$, and the SPD
matrix $\Calg_t$ that the algorithm actually applies.  They are also called
metrics because they define the norms used for uncertainty.  The reference
matrix has the following update, whereas the current and applied matrices
need not have a rank-one update:
\begin{align}
 \bVbar_1&=\lambda\bI,\nonumber\\
 \bVbar_{t+1}&=\bVbar_t+\sigma^{-2}\bq_t(a_t)\bq_t(a_t)^\top,
 \label{eq:reference-update-new}\\
 \bV_t&=\text{the exact current relinearized metric},\nonumber\\
 \Calg_t&=\text{the algorithmic SPD metric}.
\end{align}
Here $\lambda>0$ is the ridge regularization and $\sigma>0$ sets the scale
of the updates; it is the noise scale in the squared-loss specialization.
The reference update keeps each query as it was at collection time.
For the same current query, the matrices can give different widths:
\begin{equation}\label{eq:three-widths-new}
 \begin{aligned}
  \bar s_t^2(a)&=\bq_t(a)^\top\bVbar_t^{-1}\bq_t(a),\\
  s_t^2(a)&=\bq_t(a)^\top\bV_t^{-1}\bq_t(a),\\
  s_{\Calg,t}^2(a)&=\bq_t(a)^\top\Calg_t^{-1}\bq_t(a).
 \end{aligned}
\end{equation}
These are exact mathematical widths.  The computed upper width
$\hat s_t(a)$ will be a separate object, because a numerical solver need not
return the exact width of $\Calg_t$.

The frozen update also gives a way to control widths accumulated over time.
Its log-determinant growth, called the \emph{frozen information gain}, is
\begin{equation}\label{eq:frozen-gamma-new}
 \gamma_T
 :=\log\frac{\det\bVbar_{T+1}}{\det\bVbar_1}
 =\sum_{t=1}^T\log\!\left(1+\sigma^{-2}\bar s_t^2(a_t)\right).
\end{equation}
Each term records the contribution of the action actually chosen.  This
identity is why the regret proof returns to the frozen matrix even when the
policy scores actions using the current matrix.  Set $\gamma_0=0$, and define
regret as the mean reward lost relative to the best available action:
\begin{equation}\label{eq:regret-definition-new}
 \operatorname{reg}_t
 :=\sup_{a\in\cA_t}\mu_t^*(a)-\mu_t^*(a_t),
 \qquad R_T=\sum_{t=1}^T\operatorname{reg}_t.
\end{equation}

\paragraph{General protocol and timing.}
Let $(\mathsf A,\mathcal B)$ be a standard Borel action space.  At round $t$,
the learner observes a context and a nonempty random action correspondence
$\cA_t:\Omega\to 2^{\mathsf A}$ whose graph
\[
 \operatorname{Gr}(\cA_t)
 :=\{(\omega,a):a\in\cA_t(\omega)\}
 \in\cHist_t^-\otimes\mathcal B.
\]
Here $\cHist_t^-$ is the history after the context and action correspondence
are observed but before fresh round-$t$ algorithmic randomization.  The variable
$\Omega_t$ records all round-$t$ pre-reward randomization, which may be revealed
in sequential stages, and $\cHist_t=\cHist_t^-\vee\sigma(\Omega_t)$ contains
that randomization but not the reward.  The restrictions to
$\operatorname{Gr}(\cA_t)$ of
$(\omega,a)\mapsto\mu_t^*(a),\bq_t(a),m_t(a),b_t(a),\hat s_t(a)$, and, once
defined below, $U_t(a)$ are assumed jointly measurable with respect to
$\cHist_t\otimes\mathcal B$.  We also assume explicitly that
$\sup_{a\in\cA_t}\mu_t^*(a)$ and $\sup_{a\in\cA_t}U_t(a)$ are finite
$\cHist_t$-measurable random variables, that the oracle below returns a $\cHist_t$-measurable
selector in $\cA_t$, and that the simultaneous confidence and all certificate
events below, including path, operator, all-domain solver-upper, and
played-action sharpness validity, are measurable.  These are
substantive assumptions; standard-Borel structure
and graph measurability alone do not make every supremum or selector measurable.

The reference metric $\bVbar_t$ and, when used, the representation parameter
$\btheta_t$ are determined from the prefix history and are
$\cHist_t^-$-measurable.  The exact current metric $\bV_t$ is likewise formed
before the fresh certificate randomization and is $\cHist_t^-$-measurable.  The realized algorithmic operator
$\Calg_t$, path envelope $\bar D_t$, condition factors
$\kappa_{-,t},\kappa_{+,t}$, solver-width map $\hat s_t(\cdot)$, sharpness
factor $\alpha_t$, oracle tolerance $\xi_t$, and selected action $a_t$ may
depend on $\Omega_t$ and are required only to be $\cHist_t$-measurable.  Thus
``predictable'' below is reserved for $\cHist_t^-$-measurable quantities;
quantities that may use $\Omega_t$ are called pre-reward measurable.  The
reward is observed only after $a_t$ is selected.

\paragraph{What must be certified.}
The next assumption separates statistical confidence from the comparisons
needed to use it.  The confidence band concerns the center $m_t(a)$ and the
frozen width.  The path and operator certificates compare matrices.  The
solver certificate controls the width that enters the score.
Its quantifiers have different jobs: an upper-valid width is needed for every
scored action to preserve optimism, while an upper bound on excess width
(``sharpness'') is needed only for the action played to control regret.
All of these objects must be available before that action's reward is seen.

\begin{assumption}[Reference confidence and operational certificates]
\label{asm:abstract-transport}
For predictable $\beta_t\ge0$ and pre-reward measurable center $m_t(a)$ and
action-dependent envelope $b_t(a)\ge0$, a measurable simultaneous confidence
event satisfies
\begin{equation}\label{eq:abstract-confidence-new}
 |\mu_t^*(a)-m_t(a)|\le\beta_t\bar s_t(a)+b_t(a)
 \qquad\text{for every }a\in\cA_t.
\end{equation}
There is an absolutely continuous SPD path $\bV_t(\tau)$ from
$\bVbar_t$ to $\bV_t$ with a pre-reward certificate
\begin{equation}\label{eq:path-length-new}
 \begin{aligned}
  \mathcal D_t
  &:=\int_0^1\!\left\|
  \bV_t(\tau)^{-1/2}\dot\bV_t(\tau)\bV_t(\tau)^{-1/2}
  \right\|_{\rm op}d\tau\\
  &\le\bar D_t<\infty.
 \end{aligned}
\end{equation}
The algorithmic operator has a pre-reward two-sided certificate
\begin{equation}\label{eq:operator-sandwich-new}
 0<\kappa_{-,t}\le\kappa_{+,t}<\infty,\qquad
 \kappa_{-,t}\bV_t\preceq\Calg_t\preceq\kappa_{+,t}\bV_t.
\end{equation}
Finally, the solver returns a jointly measurable pre-reward width map
$\hat s_t(a)\ge0$ such that, on a measurable solver event,
\begin{equation}\label{eq:solver-interface-new}
 \begin{aligned}
  s_{\Calg,t}(a)&\le\hat s_t(a)
  &&\text{for every }a\in\cA_t,\\
  \hat s_t(a_t)&\le\alpha_t s_{\Calg,t}(a_t),
  &&\alpha_t\ge1.
 \end{aligned}
\end{equation}
The first condition is an all-domain optimism certificate.  The second is
required only at the played action.
\end{assumption}

The general score combines the center, the converted uncertainty allowance,
and the current error allowance.  The operator factor accounts for using
$\Calg_t$ in place of $\bV_t$, and $\hat s_t(a)$ is the width actually used
by the selector:
\begin{equation}\label{eq:transported-score-new}
 U_t(a)=m_t(a)
 +\beta_t e^{\bar D_t/2}\sqrt{\kappa_{+,t}}\,\hat s_t(a)
 +b_t(a),
\end{equation}
Let a pre-reward oracle return a $\cHist_t$-measurable selector
$a_t\in\cA_t$ satisfying
\begin{equation}\label{eq:approx-score-oracle-new}
 U_t(a_t)\ge\sup_{a\in\cA_t}U_t(a)-\xi_t,\qquad \xi_t\ge0.
\end{equation}
The tolerance $\xi_t$ measures how far the selected score may fall below the
best score; it is not a condition on nonlinear model training.  Regret is
defined against the mean-reward supremum, so an optimal action need not exist.

\begin{remark}[One realized solver-width map]
The all-domain solver upper condition, the played-action sharpness condition,
and the approximate-score-oracle inequality refer to the same realized width
map $\hat s_t(\cdot)$.  Refining the selected action's solve after selection
must not silently replace the value used in $U_t(a_t)$.  If refinement changes
$\hat s_t(a_t)$, the oracle inequality must be checked again, or the original
score width must remain the theorem object.  Any randomized post-selection
verification needs its own pre-draw filtration and conditional failure
allocation.  No all-action sharpness condition is required.
\end{remark}

\begin{table}[t]
\caption{Assumption ledger for abstract confidence transport.  A float64 point
residual, empirical eigenvalue estimate, or unenclosed numerical check is not a
mathematical certificate.}
\label{tab:transport-assumption-ledger}
\centering\ccetablesize
\begin{tabularx}{\columnwidth}{@{}p{.39\columnwidth}X@{}}
\toprule
Certificate & Status \\
\midrule
Reference confidence event & Statistical or model-specific premise \\
Path-length envelope & Analytic or separately certified \\
Two-sided operator comparison & Deterministic, analytic, or randomized with failure allocation \\
All-domain solver upper validity & Exact-arithmetic or verified numerical certificate \\
Played-action sharpness & Exact-arithmetic or verified selected-action certificate \\
Linearization and misspecification envelopes & Supplied pre-reward bounds \\
Measurable approximate selector & Oracle premise \\
\bottomrule
\end{tabularx}

\end{table}

\paragraph{The regret guarantee.}
Under these conditions, optimism bounds the loss at the selected action by
its uncertainty and bias terms.  Comparing that selected width back to the
reference width lets the frozen information gain control their accumulation.
The theorem states this conclusion on the joint event where the supplied
confidence and certificates hold.

\begin{theorem}[Estimator- and solver-agnostic confidence transport]
\label{thm:transported-ucb}
Under the protocol above with $\lambda>0$ and $\sigma>0$, suppose
\Cref{asm:abstract-transport} holds, the score obeys
\eqref{eq:transported-score-new}--\eqref{eq:approx-score-oracle-new}, and
$\|\bq_t(a)\|_2\le G$.  On the joint confidence and certificate event, the
instantaneous regret satisfies
\begin{equation}\label{eq:instantaneous-transport-new}
 \operatorname{reg}_t\le
 \beta_t\!\left[1+\alpha_t e^{\bar D_t}
 \sqrt{\frac{\kappa_{+,t}}{\kappa_{-,t}}}\right]\bar s_t(a_t)
 +2b_t(a_t)+\xi_t.
\end{equation}
Consequently,
\begin{align}
R_T
&\le
\sqrt{\left(\sigma^2+\frac{G^2}{\lambda}\right)\gamma_T
\sum_{t=1}^T\beta_t^2
\left[1+\alpha_t e^{\bar D_t}
\sqrt{\frac{\kappa_{+,t}}{\kappa_{-,t}}}\right]^2}
\nonumber\\[-1mm]
&\quad+2\sum_{t=1}^Tb_t(a_t)+\sum_{t=1}^T\xi_t.\label{eq:sharp-transport-regret-new}
\end{align}
In particular, the simpler bound
\begin{align}
R_T
&\le2\sqrt{\left(\sigma^2+\frac{G^2}{\lambda}\right)\gamma_T
\sum_{t=1}^T\alpha_t^2\beta_t^2e^{2\bar D_t}
\frac{\kappa_{+,t}}{\kappa_{-,t}}}
\nonumber\\[-1mm]
&\quad+2\sum_{t=1}^Tb_t(a_t)+\sum_{t=1}^T\xi_t
\label{eq:simple-transport-regret-new}
\end{align}
holds.  The bias envelope enters twice and the approximate-oracle error enters
once.  If the confidence event has probability at least $1-\delta$ and the
joint certificate event has probability at least $1-\delta_{\rm cert}$, these
bounds hold with probability at least $1-\delta-\delta_{\rm cert}$.
\end{theorem}

\paragraph{Reading the bound.}
The instantaneous bound has a statistical radius $\beta_t$, a frozen width
$\bar s_t(a_t)$, and a factor accounting for transport and computation.
The exponential $e^{\bar D_t}$ reflects two comparisons: a factor
$e^{\bar D_t/2}$ makes the current-width score optimistic, and another
compares the selected current width back to the frozen width for the regret
proof.  The ratio $\kappa_{+,t}/\kappa_{-,t}$ accounts for operator distortion,
and $\alpha_t$ accounts for excess width from the solver.  The additive terms
charge current bias and approximate score maximization.

Validity does not ensure that this bound is informative.  A large path
envelope, a weak lower operator bound, a loose solver width, or persistent
bias can make it vacuous.  In particular, an upper operator comparison alone
can support the outward comparison but cannot supply the return comparison
used by the displayed regret bound.

\paragraph{Why keep direct reference scoring as a comparator?}
Once a reference-confidence band is available, the learner can score actions
with its reference widths directly.  The next proposition states that
policy's guarantee on its own selected history.  It provides a baseline for
understanding the additional factors required by transport.

\begin{proposition}[Frozen-width UCB on its own trajectory]
\label{prop:frozen-width-ucb}
Keep the timing, measurability, bounded-query, and frozen-update assumptions
above, and assume \eqref{eq:abstract-confidence-new} along the trajectory of
the policy that uses
\[
 U_t^{\rm fr}(a)=m_t(a)+\beta_t\bar s_t(a)+b_t(a).
\]
Assume its score supremum is finite and measurable and its pre-reward
measurable selector satisfies
$U_t^{\rm fr}(a_t)\ge\sup_{a\in\cA_t}U_t^{\rm fr}(a)-\xi_t$.
On the reference-confidence event,
\begin{align}
 \operatorname{reg}_t
 &\le 2\beta_t\bar s_t(a_t)+2b_t(a_t)+\xi_t,\nonumber\\
 R_T
 &\le 2\sqrt{\left(\sigma^2+\frac{G^2}{\lambda}\right)
 \gamma_T\sum_{t=1}^T\beta_t^2}
 +2\sum_{t=1}^Tb_t(a_t)+\sum_{t=1}^T\xi_t.
 \label{eq:frozen-width-regret}
\end{align}
All quantities are evaluated on this policy's own selected history.  No path,
current-metric, or operator certificate is needed for this score.  If its
reference-confidence event has probability at least $1-\delta$, so does this
guarantee.
\end{proposition}

Holding the history fixed gives a different comparison.  If the confidence
band is the only information supplied about the mean, reducing its upper
endpoint would exclude a mean that the band still allows.  Transport cannot
justify that reduction merely by expressing the same information in another
matrix.  This is a statement about widths on the same history, not about the
regret of policies that collect different histories.

\begin{proposition}[Reference-width floor on a fixed history]
\label{prop:reference-width-floor}
On the certificate event of \Cref{asm:abstract-transport}, every scored action
satisfies
\begin{equation}\label{eq:reference-width-floor}
 w_t^{\rm tr}(a):=e^{\bar D_t/2}\sqrt{\kappa_{+,t}}\hat s_t(a)
 \ge \bar s_t(a).
\end{equation}
For exact current widths, the actionwise multiplier
\begin{equation}\label{eq:actionwise-width-ratio}
 c_t(a)=\begin{cases}
 \bar s_t(a)/s_t(a),&\bq_t(a)\ne0,\\
 1,&\bq_t(a)=0
 \end{cases}
 \qquad\text{satisfies}\qquad c_t(a)s_t(a)=\bar s_t(a).
\end{equation}
Thus this exact actionwise correction reproduces the frozen-width score when
the center, radius, and bias are the same.

More generally, fix an action and $m_t(a)$, $b_t(a)$, $\beta_t>0$, and
$\bar s_t(a)$.  If the only restriction on its mean is the full interval
\eqref{eq:abstract-confidence-new}, a score
$m_t(a)+\beta_t w_t(a)+b_t(a)$ is an upper bound for every mean in that
interval if and only if $w_t(a)\ge\bar s_t(a)$.
\end{proposition}

The last statement concerns the uncertainty set supplied by one confidence
band.  Additional restrictions on the mean, a different center, or another
simultaneous confidence event can change that set.  An exact scalar endpoint
distance supplies a worst-direction multiplier, not necessarily the
actionwise equality in \eqref{eq:actionwise-width-ratio}.  These are
same-history comparisons; they do not order the realized regret of policies
whose actions and subsequent histories differ.
Evaluating the ratio in \eqref{eq:actionwise-width-ratio} as written uses the
reference width itself and supplies no computational shortcut.  This
observation is not a general computational lower bound.

\begin{remark}[The reference band alone does not justify a smaller width]
\label{rem:naive-width-band-counterexample}
Consider an abstract confidence-band instance with $D>0$,
$\bVbar=\bI_2$, $\bV=\operatorname{diag}(e^D,e^{-D})$, and
$\bq=(1,0)^\top$.  Take $m=b=0$, $\beta>0$, and $\mu^*=\beta$.
Then $\bar s=1$ and $s=e^{-D/2}$, so
\[
 |\mu^*-m|=\beta\bar s+b,\qquad
 m+\beta s+b=\beta e^{-D/2}<\mu^*.
\]
Reference confidence therefore holds while the naive current-width score is
not an upper bound.  This illustrates the limited-information statement of
\Cref{prop:reference-width-floor}; it does not construct a realizable
scaled-tanh learning trajectory or establish that this certificate family was
necessary in the reported runs.
\end{remark}

\Cref{fig:transport-confidence-geometry} shows the illustrative slice
$D=\log 4$ and $\beta=1$.  The support of a metric unit ball in direction
$\bq$ is the corresponding inverse-metric width.  Here that support shrinks
under the current metric even though the reference confidence band still
allows its original upper endpoint.

\begin{figure}[t]
 \centering
 % Analytic illustration, not experimental data. Source: transport_theory.tex,
% rem:naive-width-band-counterexample and prop:reference-width-floor.
% Chosen slice D=log(4), beta=1, m=b=0. The metric unit balls use equal axes:
% E(Vbar) has radii (1,1); E(V) has radii (1/2,2). The plotted score bands are
% reference [-1,1], uncertified naive [-1/2,1/2], transported [-1,1].
% Native body-size TikZ; no external data, rescaling, or extra TikZ libraries.
\begin{tikzpicture}[x=1cm,y=1cm,
  every node/.style={font=\normalsize,inner sep=0pt},
  reference/.style={blue!65!black,line width=1pt},
  current/.style={orange!80!black,line width=1pt,dashed},
  transported/.style={teal!70!black,line width=1pt}]
  \path[use as bounding box] (0,-8.35) rectangle (16.2,.25);
  \node[anchor=west] at (0,0)
    {\textbf{Abstract illustration}\quad
     $\bVbar=\bI_2$, $\bV=\operatorname{diag}(e^D,e^{-D})$,
     $\bq=(1,0)^\top$, $m=b=0$.};
  \node[anchor=west] at (0,-.55)
    {Illustrative slice: $D=\log 4$, $\beta=1$; choose the allowed mean $\mu^*=\beta=1$.};

  \node[anchor=west] at (0,-1.25) {\textbf{Metric unit balls}};
  \node[anchor=west] at (0,-1.8) {$E(M)=\{z:z^\top Mz\le1\}$};
  \node[anchor=west] at (7,-1.25) {\textbf{Bands along the query direction}};

  % Equal coordinate scales: 1.05 cm per unit in both directions.
  \draw[black!40,->,line width=.45pt] (1.05,-4.35) -- (5.2,-4.35)
    node[right=3pt] {$z_1$};
  \draw[black!40,->,line width=.45pt] (3,-6.7) -- (3,-2);
  \node[anchor=west] at (3.6,-2.35) {$z_2$};
  \draw[reference,fill=blue!4] (3,-4.35) circle[radius=1.05cm];
  \draw[current,fill=orange!6,fill opacity=.55]
    (3,-4.35) ellipse[x radius=.525cm,y radius=2.1cm];
  \draw[black,->,line width=.7pt] (3,-4.35) -- (4.6,-4.35);
  \node[anchor=north] at (4.6,-4.53) {$\bq=\mathbf e_1$};
  \fill[orange!80!black] (3.525,-4.35) circle[radius=2pt];
  \fill[blue!65!black] (4.05,-4.35) circle[radius=2pt];
  \draw[orange!80!black,line width=.45pt]
    (3.525,-4.32) -- (3.525,-3.6) -- (4.4,-3.6);
  \node[anchor=west] at (4.47,-3.6) {$s=\tfrac12$};
  \draw[blue!65!black,line width=.45pt]
    (4.05,-4.32) -- (4.05,-3.02) -- (4.4,-3.02);
  \node[anchor=west] at (4.47,-3.02) {$\bar s=1$};
  \draw[reference] (0,-7.15) -- (.65,-7.15);
  \node[anchor=west] at (.82,-7.15) {Reference $E(\bVbar)$};
  \draw[current] (3.55,-7.15) -- (4.2,-7.15);
  \node[anchor=west] at (4.37,-7.15) {Current $E(\bV)$};

  % One common horizontal value scale: value v is placed at x=11.5+3v.
  \draw[black!30,line width=.45pt] (11.5,-2.1) -- (11.5,-6.35);
  \draw[black!55,densely dotted,line width=.55pt]
    (14.5,-1.95) -- (14.5,-6.35);
  \node[anchor=south] at (14.5,-1.8) {$\mu^*=1$};

  \node[anchor=west] at (7,-2.25) {Reference: $m\pm\beta\bar s$};
  \draw[reference] (8.5,-2.85) -- (14.5,-2.85);
  \draw[reference] (8.5,-2.73) -- (8.5,-2.97);
  \draw[reference] (14.5,-2.73) -- (14.5,-2.97);
  \fill[black] (14.5,-2.85) circle[radius=2.2pt];

  \node[anchor=west] at (7,-3.75)
    {Naive: $m\pm\beta s$\quad (uncertified)};
  \draw[current] (10,-4.35) -- (13,-4.35);
  \draw[current] (10,-4.23) -- (10,-4.47);
  \draw[current] (13,-4.23) -- (13,-4.47);
  \fill[black] (14.5,-4.35) circle[radius=2.2pt];
  \node[anchor=west] at (14.72,-4.35) {misses $\mu^*$};

  \node[anchor=west] at (7,-5.25)
    {Transported: $m\pm\beta e^{D/2}s$};
  \draw[transported] (8.5,-5.85) -- (14.5,-5.85);
  \draw[transported] (8.5,-5.73) -- (8.5,-5.97);
  \draw[transported] (14.5,-5.73) -- (14.5,-5.97);
  \fill[black] (14.5,-5.85) circle[radius=2.2pt];

  \draw[black!50,line width=.45pt] (8.15,-6.7) -- (14.85,-6.7);
  \foreach \x/\lab in {8.5/-1,10/-\frac12,11.5/0,13/\frac12,14.5/1}{
    \draw[black!50,line width=.45pt] (\x,-6.64) -- (\x,-6.76);
    \node[anchor=north] at (\x,-6.86) {$\lab$};
  }
  \node[anchor=north] at (11.5,-7.4) {Mean / score value};
  \node[anchor=west] at (0,-7.85)
    {$\bar s=1$, $s=e^{-D/2}$; transport restores $e^{D/2}s=\bar s$ in this direction.};
  \node[anchor=west] at (0,-8.32)
    {This is not a realized learning trajectory or an observed policy failure.};
\end{tikzpicture}

 \caption{Why the reference band alone does not justify a smaller current
 width.  In this abstract instance, the reference and current metric unit
 balls have different supports in the query direction.  With $D=\log4$,
 $m=b=0$ and $\beta=1$, the permitted mean $\mu^*=1$ exceeds the naive score
 $1/2$, whereas multiplying the current width by $e^{D/2}$ recovers the
 reference width.  The unit balls illustrate metric geometry; the current
 ball is not asserted to be a valid statistical confidence set.  This is an
 analytic illustration, not an observed failure in the scaled-tanh study.}
 \label{fig:transport-confidence-geometry}
\end{figure}

For a finite list of actions, the general selector and measurability
requirements have a concrete implementation: evaluate the scores and choose
the first maximizer in a fixed ordering.  The following corollary records
the precise version used by finite-action policies.

\begin{corollary}[Measurably enumerated finite domains]
\label{cor:finite-enumeration-new}
Suppose there is a positive $\cHist_t^-$-measurable integer $K_t$ and
$\cHist_t^-$-measurable actions $a_{t,k}:\Omega\to\mathsf A$, $k\ge1$, such
that
\[
 \cA_t=\{a_{t,1},\ldots,a_{t,K_t}\}.
\]
If the learner chooses the smallest index attaining the score maximum, then the
selected action is $\cHist_t$-measurable and
\eqref{eq:approx-score-oracle-new} holds with $\xi_t=0$.  The score and
mean-reward suprema and the universal actionwise events reduce to measurable
finite maxima and intersections.  A fixed finite action set is the immediate
special case.
\end{corollary}

\paragraph{The order of a decision round.}
Algorithm~\ref{alg:transported-ucb} uses the corrected center constructed in
\Cref{sec:corrected-center}.  The learner forms and certifies its scores
before seeing the reward, selects an action, and only then updates the frozen
statistics.  This timing keeps the collection-time queries eligible for the
sequential confidence argument.  Selection and the regret analysis must use
the same computed width map, as the remark above emphasizes.

\begin{algorithm}[t]
\caption{Corrected-center specialization of transported UCB}\label{alg:transported-ucb}
\begin{algorithmic}[1]
\STATE Observe the context and graph-measurable action correspondence $\cA_t$; this defines $\cHist_t^-$.
\STATE Take the prefix-determined $\btheta_t$, $\bVbar_t$, and $\bV_t$; using the observed context, form $\bq_t(a)$, $m_t^{\rm corr}(a)$, and $b_t(a)$.
\STATE Choose $\cHist_t^-$-measurable failure budgets before their corresponding draws; reveal the pre-reward randomness $\Omega_t$.
\STATE Construct $\Calg_t$ and obtain pre-reward certificates $\bar D_t$, $\kappa_{-,t}$, and $\kappa_{+,t}$.
\STATE Obtain one jointly measurable solver upper-width map $\hat s_t(a)$ over the score domain.
\STATE Return a $\cHist_t$-measurable $\xi_t$-approximate selector for the score in \eqref{eq:transported-score-new}.
\STATE Observe the reward; update the frozen ridge estimator and \eqref{eq:reference-update-new}.
\end{algorithmic}
\end{algorithm}

The algorithm leaves the linear algebra implementation open.  Exact Cholesky
solves satisfy the solver conditions in exact arithmetic.  Conjugate gradient
(CG) can instead solve each action system using products with a fixed SPD
operator; a residual or energy argument must certify its returned widths.
Changing the matrix through a window, subsample, or sketch is a separate
approximation and requires an operator comparison.  The controlled experiment
in \Cref{sec:transport-followup} uses dense exact-current matrices and
Cholesky solves, so it does not evaluate these approximate or matrix-free paths.

\section{How to compare the frozen and current matrices}
\label{sec:logarithmic-transport}

The main theorem assumes a comparison between matrices.  This section explains
how to obtain one when past gradients are recomputed at the current parameter.
The matrix order $A\preceq B$ means that $B-A$ is positive semidefinite: the
quadratic form of $B$ is at least that of $A$ in every direction.  A two-sided
comparison therefore controls all query widths at once.

\paragraph{Measure relative change along a path.}
Imagine continuously changing the frozen matrix into the current matrix.
The path-length integrand below measures the largest relative change at each
point, after normalizing by the matrix there.  Integrating those changes gives
a multiplicative comparison between the endpoints.  The exponential form
allows large changes; it does not promise that the resulting factor is small.

\begin{lemma}[Logarithmic comparison along an SPD path]
\label{lem:log-spd-path}
Let $\bV:[0,1]\to\mathbb{S}_{++}^d$ be entrywise absolutely continuous and
let $\dot\bV(\tau)$ denote its almost-everywhere derivative.  If
\begin{equation}\label{eq:generic-path-length-new}
 \mathcal D(\bV)=\int_0^1
 \|\bV(\tau)^{-1/2}\dot\bV(\tau)\bV(\tau)^{-1/2}\|_{\rm op}\,d\tau,
\end{equation}
then
\begin{equation}\label{eq:log-path-sandwich-new}
 e^{-\mathcal D(\bV)}\bV(0)\preceq\bV(1)
 \preceq e^{\mathcal D(\bV)}\bV(0).
\end{equation}
No commutativity or small-path assumption is required.
\end{lemma}

\begin{remark}[Thompson distance and selected-path length]
The integrand in \eqref{eq:generic-path-length-new} is the classical Thompson
Finsler norm on the positive-definite cone
\citep{thompson1963certain,nussbaum1994finsler}.  Its endpoint metric is
\[
 d_{\rm Th}(A,B)
 :=\inf\{r\ge0:e^{-r}A\preceq B\preceq e^rA\}.
\]
Equivalently,
\[
 d_{\rm Th}(A,B)
 =\left\|\log(A^{-1/2}BA^{-1/2})\right\|_{\rm op}.
\]
Therefore \Cref{lem:log-spd-path} gives
\begin{align*}
 d_{\rm Th}(\bV(0),\bV(1))&\le\mathcal D(\bV),\\
 d_{\rm Th}(\bVbar_t,\bV_t)&\le\mathcal D_t\le\bar D_t.
\end{align*}
The endpoint distance depends only on the two metrics, whereas
$\mathcal D(\bV)$ depends on the selected path.  The direct SPD path length is
independent of a factorization; the upper bound from
\Cref{lem:factor-spd-path} depends on the chosen factorization and interpolation.
\end{remark}

\begin{remark}[Sharpness of the symmetric exponential factors]
For $D\ge0$, let
\[
\begin{aligned}
 \bVbar&=\bI_2,&
 \bV&=\begin{pmatrix}e^D&0\\0&e^{-D}\end{pmatrix},\\
 \bV(\tau)&=\begin{pmatrix}e^{D\tau}&0\\0&e^{-D\tau}\end{pmatrix}.
\end{aligned}
\]
Then
\[
 \bV(\tau)^{-1/2}\dot\bV(\tau)\bV(\tau)^{-1/2}
 =\begin{pmatrix}D&0\\0&-D\end{pmatrix},
\]
so $\mathcal D(\bV)=d_{\rm Th}(\bVbar,\bV)=D$.  For $q_1=e_1$,
$\bar s(q_1)=1$ and $s(q_1)=e^{-D/2}$, so
$\bar s=e^{D/2}s$ is attained.  For $q_2=e_2$,
$s(q_2)=e^{D/2}$ and $\bar s(q_2)=1$, so
$s=e^{D/2}\bar s$ is attained.  Thus each one-way factor is sharp, and a
uniform argument based only on the symmetric Thompson-distance bound cannot
improve the resulting round-trip factor $e^D$.  This is a metric-factor
sharpness statement, not a contextual-bandit regret lower bound and not a
claim that every trajectory attains the worst case.
\end{remark}

\paragraph{Work with the gradients that form the matrix.}
A replay matrix is a ridge term plus a Gram matrix of gradient features.
Stack those features as rows of $\Bmat$.  It is often easier to bound how
these rows change than to analyze the full SPD path directly.  The next lemma
turns a bound on the changing factor into a bound on matrix path length.
Its normalization by $\bV(\tau)^{-1/2}$ retains information about which
directions the data have already constrained.

\begin{lemma}[Rectangular factor-path certificate]
\label{lem:factor-spd-path}
Let $\Bmat:[0,1]\to\R^{m\times d}$ be absolutely continuous and set
$\bV(\tau)=\lambda\bI+\Bmat(\tau)^\top\Bmat(\tau)$, with $\lambda>0$.
Then
\begin{equation}\label{eq:factor-path-bound-new}
 \mathcal D(\bV)\le2\int_0^1
 \|\dot\Bmat(\tau)\bV(\tau)^{-1/2}\|_{\rm op}\,d\tau.
\end{equation}
More generally, suppose there is a measurable function
$\nu_t:[0,1]\to[0,\infty)$ with $\nu_t\in L^1([0,1])$ such that
\[
 \dot\Bmat(\tau)^\top\dot\Bmat(\tau)
 \preceq\nu_t(\tau)^2\bV(\tau)
\]
almost everywhere.  Then
$\mathcal D(\bV)\le2\int_0^1\nu_t(\tau)d\tau$.  For nonnegative $\nu_t$, the
displayed PSD premise is equivalently the norm condition
$\|\dot\Bmat(\tau)\bV(\tau)^{-1/2}\|_{\rm op}\le\nu_t(\tau)$ almost everywhere.
\end{lemma}

For a scalar prediction, each past observation contributes a gradient row.
Move that row from its collection-time parameter $\btheta_s$ to the current
parameter $\btheta_t$ along its own straight segment.  The $s$th row is
\begin{equation}\label{eq:scalar-factor-path-new}
 \sigma^{-1}\nabla_\theta\mu_{\btheta_s+
 \tau(\btheta_t-\btheta_s)}(x_s,a_s)^\top.
\end{equation}
The endpoint matrices are the frozen and exact current matrices.  The path
between them is a device for comparing those endpoints: different rows begin
at different collection-time parameters.  It need not be the generalized
Gauss--Newton (GGN) matrix of any common intermediate model, and it is not the
learner's training trajectory.
A weighted or heteroscedastic squared loss uses row
$\sqrt{w_s}\,\nabla\mu^\top$ for a nonnegative collection-time weight $w_s$
fixed before the corresponding reward.
If the weight is fixed along the interpolation, the same proof only rescales
the corresponding row.  If it changes with the path, the derivative of
$\sqrt{w_s(\tau)}$ is an additional obligation.

\paragraph{A certificate from smoothness and parameter movement.}
A prediction-Hessian bound controls how fast each gradient can change along
its segment.  Combined with the ridge lower bound on the matrix, it gives a
certificate in terms of $Q_t$, the accumulated squared distances from the
current parameter to collection-time parameters.  This route is easy to state,
but it can discard information about well-constrained directions.

\begin{corollary}[Hessian path certificate]
\label{cor:hessian-path-new}
Suppose the parameter region is convex and
$\|\nabla_\theta^2\mu_\theta(x_s,a_s)\|_{\rm op}\le L_g$ along every segment
in \eqref{eq:scalar-factor-path-new}.  With
\begin{equation}\label{eq:Qt-new}
 Q_t=\sum_{s<t}\|\btheta_t-\btheta_s\|_2^2,
\end{equation}
the exact current and frozen metrics satisfy
\begin{equation}\label{eq:hessian-path-sandwich-new}
 e^{-\bar D_t}\bVbar_t\preceq\bV_t\preceq e^{\bar D_t}\bVbar_t,
 \qquad
 \bar D_t=\frac{2L_g\sqrt{Q_t}}{\sigma\sqrt\lambda}.
\end{equation}
This remains valid when $\bar D_t\ge1$.  At $t=1$, the factor has zero rows,
$Q_1=0$, and both metrics equal $\lambda\bI$.
\end{corollary}

The relative condition in \Cref{lem:factor-spd-path} retains that directional
information and can be much sharper than the ridge-only Hessian bound.
Changes in directions with larger matrix eigenvalues are downweighted by
the inverse-matrix normalization.  Both constructions permit gradient
rotation; the factor-path lemma requires neither a fixed subspace nor known
rank.

\paragraph{Convert matrix comparisons into width comparisons.}
The next proposition gives the outward and return comparisons used in the
regret theorem.  It also shows how a relative-error bound for an approximate
operator supplies the factors $\kappa_{-,t}$ and $\kappa_{+,t}$.

\begin{proposition}[Width transport and approximate operators]
\label{prop:width-transport-new}
Under \eqref{eq:path-length-new} and \eqref{eq:operator-sandwich-new},
\begin{align}
 \bar s_t^2(a)&\le e^{\bar D_t}\kappa_{+,t}s_{\Calg,t}^2(a),
 &s_{\Calg,t}^2(a)&\le
 e^{\bar D_t}\kappa_{-,t}^{-1}\bar s_t^2(a).
 \label{eq:two-way-width-new}
\end{align}
Exact current curvature has $\kappa_{-,t}=\kappa_{+,t}=1$.  More generally,
the operational relative-error certificate
\begin{equation}\label{eq:relative-operator-certificate-new}
 \|\bV_t^{-1/2}(\Calg_t-\bV_t)\bV_t^{-1/2}\|_{\rm op}
 \le\varepsilon_t<1
\end{equation}
gives $\kappa_{-,t}=1-\varepsilon_t$ and
$\kappa_{+,t}=1+\varepsilon_t$.  An upper comparison alone does not imply a
lower factor.
\end{proposition}

\paragraph{Compare the endpoints directly when possible.}
Following an entire path is not required if a valid two-sided comparison of
the endpoint matrices is already available.  The next corollary inserts that
comparison into the same confidence and regret argument.  Its upper factor
protects optimism; its lower factor supplies the return comparison, so their
ratio enters the regret bound.

\begin{corollary}[Direct endpoint sandwich]
\label{cor:relative-endpoint-ucb}
Replace the path premise \eqref{eq:path-length-new} in
\Cref{thm:transported-ucb} by pre-reward measurable factors that, on the joint
certificate event, satisfy
\begin{equation}\label{eq:relative-endpoint-sandwich}
 0<c_{-,t}\le c_{+,t}<\infty,\qquad
 c_{-,t}\bVbar_t\preceq\bV_t\preceq c_{+,t}\bVbar_t.
\end{equation}
Replace \eqref{eq:transported-score-new} by the score
\begin{equation}\label{eq:relative-endpoint-score}
 U_t^{\rm rel}(a)=m_t(a)+\beta_t\sqrt{c_{+,t}\kappa_{+,t}}
 \hat s_t(a)+b_t(a),
\end{equation}
and require the same finite measurable supremum and approximate-selector
condition for this score.  Keep every other hypothesis of that theorem.
On the joint confidence and certificate event, writing
$h_t=1+\alpha_t\sqrt{c_{+,t}\kappa_{+,t}/(c_{-,t}\kappa_{-,t})}$,
\begin{align}
 \operatorname{reg}_t
 &\le\beta_t h_t\bar s_t(a_t)+2b_t(a_t)+\xi_t,\nonumber\\
 R_T&\le\sqrt{\left(\sigma^2+\frac{G^2}{\lambda}\right)
 \gamma_T\sum_{t=1}^T\beta_t^2h_t^2}
 +2\sum_{t=1}^Tb_t(a_t)+\sum_{t=1}^T\xi_t.
 \label{eq:relative-endpoint-regret}
\end{align}
The probability statement of \Cref{thm:transported-ucb} is unchanged, with
the endpoint certificate included in its joint certificate event.
\end{corollary}

This result uses the endpoint matrices directly.  In particular,
$\max\{-\log c_{-,t},\log c_{+,t}\}$ bounds their Thompson distance;
it need not bound the length of the synthetic path
\eqref{eq:scalar-factor-path-new}.  Its effective width still obeys the
reference-width floor: $\sqrt{c_{+,t}\kappa_{+,t}}\hat s_t(a)\ge\bar s_t(a)$.

\paragraph{What extends to vector outputs?}
The factor construction also accommodates several output coordinates and
weights, provided their changes are controlled as well.  For vector outputs,
let $J_s(\theta)\in\R^{p_s\times d}$ and supply an
absolutely continuous factor $R_s(\theta)\in\R^{r_s\times p_s}$ satisfying
$R_s^\top R_s=W_s\succeq0$.  Stacking
$B_s=R_sJ_s\in\R^{r_s\times d}$ gives
$\lambda I+\sum_sJ_s^\top W_sJ_s$.  Along a path,
$\dot B_s=\dot R_sJ_s+R_s\dot J_s$; both terms must be controlled.  This
covers the geometry of weighted Fisher/GGN factors but does not by itself prove
a likelihood confidence sequence or an exponential-family regret theorem.

\section{A confidence center that does not require optimizer convergence}
\label{sec:corrected-center}

So far, the confidence band has been an assumption.  We now construct one
for squared loss.  The difficulty is that an arbitrarily trained nonlinear
predictor need not be the center of a valid confidence band.  The construction
uses that predictor as an anchor, then corrects it using a linear estimate
formed from collection-time features.  Confidence comes from this frozen
linear estimation problem; it does not require the nonlinear optimizer to
have converged.

There is a cost to this separation.  Linearizing the model introduces Taylor
error, and the reference model may not exactly represent the true mean.
The resulting band must account for both errors, in past observations and in
the current prediction.  The assumptions below make these costs explicit.

\paragraph{Noise and the reference model.}
Let $\sigma>0$ and let the reward obey
$r_t=\mu_t^*(a_t)+\eta_t$, where, for every $u\in\R$,
\begin{equation}\label{eq:new-noise-condition}
 \mathbb E[\exp(u\eta_t)\mid\cHist_t]
 \le\exp(u^2\sigma^2/2).
\end{equation}
Thus the noise condition is imposed after the action is selected and before
the reward is revealed.  The ridge penalty is centered at zero in the chosen
parameter coordinates, and every norm bound below, including
$\|\btheta^\circ\|_2\le S$, uses those coordinates.  For a pretrained
initialization $\btheta_0$, the intended convention is displacement coordinates
$\btheta=\btheta_{\rm orig}-\btheta_0$; then $S$ bounds the displacement of the
reference model from initialization rather than its raw parameter norm.  Fix
before data collection a reference parameter $\btheta^\circ\in\Theta$ with
$\|\btheta^\circ\|_2\le S$, and allow
approximate realizability
\begin{equation}\label{eq:approx-realizability-new}
 \mu^*(z)=\mu_{\btheta^\circ}(z)+m^\circ(z),\qquad
 |m^\circ(z)|\le\epsilon_{\rm mis}(z).
\end{equation}
\paragraph{Turn past rewards into a frozen linear estimation problem.}
The predictable sequence $\btheta_t\in\Theta$ is otherwise arbitrary.  Set
$\bq_t(a)=\nabla_\theta\mu_{\btheta_t}(x_t,a)$.  A collection-time
linearization has a known nonlinear offset and a linear term in the reference
parameter.  Subtracting that offset from the reward gives a
\emph{pseudo-response}: an approximate linear observation of the fixed
reference parameter, with noise and Taylor and misspecification errors.
The pseudo-responses and their frozen ridge estimate are
\begin{align}
 y_s&=r_s-\mu_{\btheta_s}(x_s,a_s)+\bq_s(a_s)^\top\btheta_s,
 \label{eq:pseudo-response-new}\\
 \hat\btheta_t^{\rm lin}
 &=\bVbar_t^{-1}\sigma^{-2}\sum_{s<t}\bq_s(a_s)y_s.
 \label{eq:frozen-ridge-new}
\end{align}
\paragraph{Correct the current prediction.}
The linear estimate uses old, predictable features, while the action query
uses the current model.  The corrected center connects them:
\begin{equation}\label{eq:corrected-center-new}
 m_t^{\rm corr}(a)=\mu_{\btheta_t}(x_t,a)
 +\bq_t(a)^\top(\hat\btheta_t^{\rm lin}-\btheta_t).
\end{equation}

The correction applies the current gradient to the difference between the
frozen linear estimate and the current parameter.  Thus the certified center
is generally different from the raw prediction of the trained model.  The
identity supporting its confidence band holds for any predictable parameter
path in the certified region, but it does not make the approximation errors
small.

\paragraph{Account for historical and current errors separately.}
Fix $\delta\in(0,1)$.  Let $\rho_s^\circ$ be the Taylor remainder of
$\mu_{\btheta^\circ}$ around $\btheta_s$, and supply envelopes
\begin{equation}\label{eq:lin-mis-envelopes-new}
 |\rho_s^\circ|\le\epsilon_{\rm lin}(s),\qquad
 |m^\circ(x_s,a_s)|\le\epsilon_{\rm mis}(s),
\end{equation}
with the corresponding actionwise current envelopes
$\epsilon_{\rm lin,t}(a)$ and $\epsilon_{\rm mis,t}(a)$.  Each envelope is
deterministic or measurable before the reward at the round where it is used;
the complete historical sum below is known before selecting $a_t$.
Past errors affect the fitted linear estimate and therefore enlarge the
confidence radius.  Current errors enter the final prediction band
additively.  Define
\begin{equation}\label{eq:beta-corrected-new}
 \begin{aligned}
  \beta_t^{\rm corr}
  &=\sqrt{\gamma_{t-1}+2\log(1/\delta)}+\sqrt\lambda S\\
  &\quad+\frac1\sigma\sqrt{\sum_{s<t}
  [\epsilon_{\rm lin}(s)+\epsilon_{\rm mis}(s)]^2}.
 \end{aligned}
\end{equation}

The radius contains the stochastic uncertainty term, the ridge penalty for
the reference parameter, and the accumulated historical approximation error.
The following result supplies the reference-confidence band required by
\Cref{asm:abstract-transport}.

\begin{corollary}[Corrected-center confidence without optimizer assumptions]
\label{cor:corrected-center-new}
With probability at least $1-\delta$, simultaneously over rounds and actions,
\begin{equation}\label{eq:corrected-confidence-new}
 |\mu_t^*(a)-m_t^{\rm corr}(a)|
 \le\beta_t^{\rm corr}\bar s_t(a)
 +\epsilon_{\rm lin,t}(a)+\epsilon_{\rm mis,t}(a).
\end{equation}
The identity and bound hold for any predictable $\btheta_t$ in the certified
region.  They require no exact ERM, approximate stationarity, optimizer
residual, or optimizer convergence rate.
\end{corollary}

The historical deterministic term in \eqref{eq:beta-corrected-new} follows from
$\Phi_t^\top\bVbar_t^{-1}\Phi_t\preceq I$ for
$\Phi_t=[\bq_1/\sigma,\ldots,\bq_{t-1}/\sigma]$.  Its scaling is
$1/\sigma$, not $1/\sigma^2$.  Fixed nonvanishing misspecification generally
creates a linear current-bias term and can also make the historical radius
vacuous.  This is a graceful-degradation statement, not a robust
misspecified-bandit guarantee.

\section{When the regret bound is informative}
\label{sec:closed-rates-new}

The general theorem depends on quantities accumulated along the policy's
history.  This section replaces the information gain by a bound in terms of
parameter dimension and horizon, then asks when the remaining transport and
approximation terms are controlled.  A valid certificate alone is not enough:
a large historical error radius or persistent current bias can still prevent
the displayed regret bound from being sublinear.

\paragraph{Bound information gain using dimension.}
Let $r_T=\min\{d,T\}$.  The frozen Gram matrix has at most this rank above
its ridge term.  Bounding its eigenvalue sum using the query-norm bound gives
\begin{equation}\label{eq:finite-d-gamma-new}
 \gamma_T\le\Gamma_{T,d}:=
 r_T\log\!\left(1+\frac{TG^2}{r_T\lambda\sigma^2}\right).
\end{equation}
No fixed tangent subspace is required.  The next corollary substitutes this
bound and the corrected-center radius into the general theorem.  Its
$B_T$ collects historical error energy, while $E_T^{\rm cur}$ collects the
current prediction-error allowances.

\begin{corollary}[Finite-dimensional corrected-center rate]
\label{cor:finite-d-rate-new}
Suppose $\bar D_t\le D$, $\alpha_t\le\alpha$, and
$\kappa_{+,t}/\kappa_{-,t}\le K$ for all $t$.  Let
\begin{align*}
 B_T&=\sum_{s<T}[\epsilon_{\rm lin}(s)+\epsilon_{\rm mis}(s)]^2,\\
 E_T^{\rm cur}&=\sum_{t=1}^T
 [\epsilon_{\rm lin,t}(a_t)+\epsilon_{\rm mis,t}(a_t)].
\end{align*}
Then \Cref{thm:transported-ucb,cor:corrected-center-new} imply
\begin{align}
R_T
&\le2\alpha e^D\sqrt K\nonumber\\
&\quad{}\times
\left[\sqrt{\Gamma_{T,d}+2\log(1/\delta)}+\sqrt\lambda S
+\frac{\sqrt{B_T}}{\sigma}\right]\nonumber\\
&\quad{}\times
\sqrt{\left(\sigma^2+\frac{G^2}{\lambda}\right)T\Gamma_{T,d}}
\nonumber\\
&\quad+2E_T^{\rm cur}+\sum_{t=1}^T\xi_t.
\label{eq:finite-d-rate-new-display}
\end{align}
With exact realizability and a frozen linear model, $B_T=E_T^{\rm cur}=0$.
For fixed $d$ and problem constants, the leading term is then
$O(e^D\sqrt K\,d\sqrt T\log T)$.  Exact solves set $\alpha=1$, exact current
curvature sets $K=1$, and exact score maximization sets $\xi_t=0$, recovering
the standard LinUCB determinant-potential rate up to displayed constants.
\end{corollary}

\paragraph{Control both transport and linearization error.}
The preceding bound becomes useful when the transport factors stay controlled
and the approximation errors are small enough.  The next corollary gives a
sufficient regime: the model is nearly linear throughout a certified parameter
region.  Its scale $W$ controls both gradient change and Taylor error.  The
condition relates this scale to the horizon; it is not a statement that an
arbitrary fixed nonlinear model has sublinear regret.

\begin{corollary}[Global near-linearity under exponential transport]
\label{cor:global-near-linear-new}
Assume exact realizability, $\mu^*(z)=\mu_{\btheta^\circ}(z)$, and suppose
$\|\btheta^\circ\|_2\le R$ and $\|\btheta_t\|_2\le R$ for every $t$.
On the certified convex radius-$R$ region, assume the prediction gradient is
$L_\mu$-Lipschitz for the Taylor envelopes and satisfies the replay-path
Hessian bound $L_g$ in \Cref{cor:hessian-path-new}, with
\[
 L_\mu\le\frac{c_\mu}{\sqrt W},\qquad
 L_g\le\frac{c_g}{\sqrt W},
\]
with $W>0$.  Use the corrected center, exact current curvature, exact dense
solves, and exact score maximization in finite parameter dimension $d$, and
suppose $\|\bq_t(a)\|_2\le G$.  Thus
$\alpha_t=1$, $\kappa_{-,t}=\kappa_{+,t}=1$, and $\xi_t=0$.
Define
\[
 D_{T,W}:=\frac{4c_gR\sqrt{T-1}}{\sigma\sqrt{\lambda W}}.
\]
Then, with probability at least $1-\delta$,
\begin{align}
 R_T
 &\le 2e^{D_{T,W}}
 \Biggl[
 \begin{aligned}
  &\sqrt{\Gamma_{T,d}+2\log(1/\delta)}+\sqrt\lambda R\\[-1mm]
  &\quad+\frac{2c_\mu R^2\sqrt T}{\sigma\sqrt W}
 \end{aligned}
 \Biggr]\nonumber\\[-1mm]
 &\quad\times
 \sqrt{\sigma^2+\frac{G^2}{\lambda}}\,
 \sqrt{T\Gamma_{T,d}}\nonumber\\
 &\quad+\frac{4c_\mu R^2T}{\sqrt W}.
 \label{eq:global-near-linear-new}
\end{align}
Consequently, $W=\Omega(T)$ with a sufficiently large problem-dependent
constant keeps $D_{T,W}=O(1)$.  For fixed $d$ and fixed problem constants,
\[
 R_T=O(d\sqrt T\log T)+O(\sqrt T).
\]
This conclusion requires no $D_t<1$ condition.  Here $W$ is an abstract
near-linearity scale, not generic network width or parameter count, and the
result is not an unrestricted neural-training theorem.  No optimizer-residual
rate is required because the corrected center is used.
\end{corollary}

\begin{remark}[Comparable guarantees do not order policies]
\Cref{prop:frozen-width-ucb,cor:corrected-center-new} give a frozen-width
counterpart to \eqref{eq:finite-d-rate-new-display}, with the factor
$\alpha e^D\sqrt K$ removed and all remaining quantities evaluated on the
frozen policy's own history.  With the same $\lambda$, $\sigma$, and
confidence level, common uniform bounds on the query norm, reference norm,
approximation envelopes, and accumulated oracle error give guarantees of the
same order when $D$, $\alpha$, and $K$ are bounded.  In particular, under the
statistical, uniform-region, and smoothness assumptions of
\Cref{cor:global-near-linear-new}, the corrected center and exact
frozen-width score maximization give the same
$O(d\sqrt T\log T)+O(\sqrt T)$ rate when $W=\Omega(T)$, for fixed $d$
and fixed problem constants.
This comparison does not rank realized regret: different selected actions
change the representation paths, centers, radii, and information gains.
It neither proves that transport is unnecessary nor rules out improved
performance from current curvature.  The stated $W$ scaling is a sufficient
condition for these bounds, not a necessary condition for an informative
transport certificate.
\end{remark}

\paragraph{A concrete model with fixed feature directions.}
The scaled-tanh model makes the smoothness condition explicit.  It also has
more structure than the general Hessian argument uses: changing the parameter
rescales each feature gradient without rotating its direction.  This permits
the sharper endpoint comparison developed after the model's derivative
formulas.

\begin{remark}[Scaled-tanh specialization]
For
\[
 \mu_{\btheta,W}(z)=\sqrt W\tanh\!\left(
 \frac{\phi(z)^\top\btheta}{\sqrt W}\right),
 \qquad \|\phi(z)\|_2\le B_\phi,
\]
write $u=\phi(z)^\top\btheta/\sqrt W$.  Then
\begin{align*}
 \nabla_\theta\mu_{\btheta,W}(z)
 &=\operatorname{sech}^2(u)\phi(z),\\
 \nabla_\theta^2\mu_{\btheta,W}(z)
 &=-\frac{2}{\sqrt W}\operatorname{sech}^2(u)\tanh(u)
 \phi(z)\phi(z)^\top.
\end{align*}
Since
$\sup_u|2\operatorname{sech}^2(u)\tanh(u)|=4/(3\sqrt3)$,
\Cref{cor:global-near-linear-new} applies with
\[
 G=B_\phi,\qquad
 c_\mu=c_g=c_h:=\frac{4B_\phi^2}{3\sqrt3}.
\]
The symbol $W$ is a smoothness scale here; it is not identified with hidden
width or parameter count.
\end{remark}

\paragraph{Compare scalar weights instead of accumulating a path bound.}
For a fixed observation, the gradients at collection time and at the current
parameter point in the same feature direction.  Their contributions to the
Gram matrix differ only by the \emph{squares} of their scalar gradient
multipliers.  Bounds on the ratio of these weights therefore compare the
endpoint matrices directly.  The common ridge term must be included in that
comparison as well.

There are two versions.  A uniform bound uses only the parameter and feature
norm bounds.  When observations share a finite collection of feature vectors,
a category-wise bound uses the accumulated collection-time weight and the
current weight within each category.  Both are certificates for the same
reference-confidence construction; neither supplies a new statistical event.

\begin{corollary}[Relative weights for scaled tanh]
\label{cor:tanh-relative-weights}
For the scaled-tanh model above, suppose $W>0$, $\|\phi(z_s)\|_2\le B_\phi$,
and $\|\btheta_s\|_2,\|\btheta_t\|_2\le R$ for all $s<t$.
Set $\omega_\theta(z)=\operatorname{sech}^2(\phi(z)^\top\theta/\sqrt W)$ and
$\ell=\operatorname{sech}^2(B_\phi R/\sqrt W)>0$.  For
\begin{align*}
 \bVbar_t&=\lambda\bI+\sigma^{-2}\sum_{s<t}
 \omega_{\btheta_s}(z_s)^2\phi(z_s)\phi(z_s)^\top,\\
 \bV_t&=\lambda\bI+\sigma^{-2}\sum_{s<t}
 \omega_{\btheta_t}(z_s)^2\phi(z_s)\phi(z_s)^\top,
\end{align*}
the uniform endpoint certificate is
\begin{equation}\label{eq:tanh-uniform-relative}
 \ell^2\bVbar_t\preceq\bV_t\preceq\ell^{-2}\bVbar_t,
 \qquad d_{\rm Th}(\bVbar_t,\bV_t)
 \le4\log\cosh(B_\phi R/\sqrt W).
\end{equation}
If the observed feature vectors belong to a finite set $\{\phi_c\}$, let
$n_{t,c}$ be their counts, let
$f_{t,c}=\sum_{s<t:\phi(z_s)=\phi_c}\omega_{\btheta_s}(z_s)^2$, and put
$g_{t,c}=n_{t,c}\operatorname{sech}^4(\phi_c^\top\btheta_t/\sqrt W)$.
The sharper category certificate is
\begin{equation}\label{eq:tanh-category-relative}
 c_{-,t}=\min\!\left\{1,\min_{c:n_{t,c}>0}\frac{g_{t,c}}{f_{t,c}}\right\},
 \qquad
 c_{+,t}=\max\!\left\{1,\max_{c:n_{t,c}>0}\frac{g_{t,c}}{f_{t,c}}\right\},
\end{equation}
with $c_{-,1}=c_{+,1}=1$ for an empty history.  It satisfies
\eqref{eq:relative-endpoint-sandwich} and
$\ell^2\le c_{-,t}\le1\le c_{+,t}\le\ell^{-2}$.
Either certificate can be used in \Cref{cor:relative-endpoint-ucb}.
\end{corollary}

The uniform bound has no explicit horizon dependence when its norm bounds and
$W$ are fixed.  This controls transport, but does not remove the separate
historical and current approximation errors from the regret bound.
For the category version, maintaining a count and a collection-weight sum
per category lets \eqref{eq:tanh-category-relative} use current category
projections without an eigensolve or a scan of the full history.  It reduces
the work needed for this certificate, not the cost of all width or center
solves.  The construction relies on fixed gradient directions and does not
cover general gradient rotation or establish improved policy regret.
The appendix gives the termwise proof and an example where the endpoint
distance is zero although the synthetic path has positive length.

\paragraph{A local comparison when changes are small.}
A local two-sided comparison gives another way to relate the same matrices.
It replaces the exponential factors by a local sandwich that requires a
smallness condition.
If $\chi_t<1$ and
\begin{equation}\label{eq:local-additive-new}
 (1-\chi_t)^2\bVbar_t\preceq\bV_t
 \preceq(1+\chi_t)^2\bVbar_t,
\end{equation}
use the local score obtained from \eqref{eq:transported-score-new} by replacing
$e^{\bar D_t/2}$ with $1+\chi_t$.
\begin{corollary}[Local small-drift transport]
\label{cor:local-small-drift-new}
Under the remaining hypotheses of \Cref{thm:transported-ucb}, the local score
satisfies
\begin{align}
R_T
&\le\sqrt{\left(\sigma^2+\frac{G^2}{\lambda}\right)\gamma_T
\sum_{t=1}^T\beta_t^2
\left[1+\alpha_t\frac{1+\chi_t}{1-\chi_t}
\sqrt{\frac{\kappa_{+,t}}{\kappa_{-,t}}}\right]^2}
\nonumber\\
&\quad+2\sum_{t=1}^Tb_t(a_t)+\sum_{t=1}^T\xi_t.
\label{eq:local-small-drift-regret-new}
\end{align}
\end{corollary}
Neither the local factor nor the exponential factor uniformly dominates the
other because they summarize different paths.

\begin{table}[t]
\caption{Transport certificates used in this paper.  The two-way factor is
the product of the outward and return width factors.  ``Smallness'' states
the additional restriction needed for that factor; ``rotation'' records
whether replay gradients may change direction. The final two rows give only
the operator-comparison factor; a separate metric-transport factor is still
required.}
\label{tab:transport-comparison-new}
\centering\ccetablesize
\begin{tabularx}{\columnwidth}{@{}
  >{\raggedright\arraybackslash}p{.17\columnwidth}
  >{\raggedright\arraybackslash}X
  >{\raggedright\arraybackslash}p{.18\columnwidth}
  >{\centering\arraybackslash}p{.10\columnwidth}
  >{\centering\arraybackslash}p{.10\columnwidth}@{}}
\toprule
Route & Assumption & two-way factor & smallness & rotation \\
\midrule
Logarithmic path & normalized SPD path length $\bar D_t$ & $e^{\bar D_t}$ & no & yes \\
Direct endpoint & \eqref{eq:relative-endpoint-sandwich} & $\sqrt{c_+/c_-}$ & no & yes \\
Category weights & fixed feature directions; \eqref{eq:tanh-category-relative} & $\sqrt{c_+/c_-}$ & no & no \\
Local sandwich & \eqref{eq:local-additive-new} & $(1+\chi_t)/(1-\chi_t)$ & $\chi_t<1$ & yes \\
Exact current & $\Calg_t=\bV_t$ & $1$ & no & yes \\
Approximate current & two-sided $\kappa_\pm$ certificate & $\sqrt{\kappa_+/\kappa_-}$ & $\kappa_->0$ & yes \\
\bottomrule
\end{tabularx}

\end{table}

\paragraph{What remains model-specific.}
The geometric comparisons can be reused with another estimator if it supplies
\eqref{eq:abstract-confidence-new}.  Geometry alone does not supply that
statistical premise.  For an exponential-family model, one would also need a
likelihood-score confidence event, a conversion to prediction error,
pre-reward collection-time Fisher weights, control of weight and Jacobian
transport, and an information-potential bound.  Those additional steps are
outside the squared-loss specialization established here.
